\thispagestyle{empty}\vspace*{\fill}
\begin{minipage}{128mm}\fontsize{9}{13}\selectfont\color{otMuted}
{\otSemi\color{otInk}Mathématiques · Sixième · France}\par\medskip
Conception et mise en page : collection OptiTAB.\par
315 exercices de chapitre, 21 défis et deux bilans originaux, tous corrigés.\par
Édition du 23 octobre 2026 · Année scolaire 2026-2027.\par\bigskip
Cette édition suit le programme de mathématiques du cycle 3 publié au
BO n° 16 du 17 avril 2025, applicable en Sixième depuis la rentrée 2025.
Les rubriques « Sixième » de chaque domaine ont guidé le périmètre.\par\bigskip
L’ordre des chapitres, les situations et les personnages sont des choix
éditoriaux. Les défis sont facultatifs. Les bilans sont des entraînements
originaux, sans statut d’examen officiel.\par\bigskip
Les données numériques des problèmes sont construites pour apprendre ;
les données réelles citées (population mondiale, distances dans le
système solaire, calendriers) sont arrondies. Les figures de principe ne
sont pas à mesurer, sauf consigne explicite ; graduations et quadrillages
font foi.\par\bigskip
Les réponses ajoutées en vert aux figures corrigées ne sont pas des données
de l’énoncé. Les calculs du livre se font sans calculatrice, sauf mention
contraire dans un défi.\par\bigskip
© OptiTAB, 2026.\par\otSite
\end{minipage}\clearpage
